\section{Simulation-to-Reality Analysis}
\label{sec:simreal}

\subsection{Where the two domains can be compared at all}

The command stream is recorded on the same topic,
\texttt{/planner/cmd\_vel\_safe}, at the same boundary in both domains,
and is therefore the one axis on which a simulated and a physical run
measure the same signal with the same instrument. Everything spatial is
excluded by construction (Section~\ref{sec:realresults}).

What follows is not a test of predictive validity. The frozen and the
flown stack differ in the eleven ways of Table~\ref{tab:gap}, so
agreement would not establish that the calibration worked and
disagreement would not establish that it failed. The question is
narrower: how far apart are the two domains where both measure the same
thing, and does the distance move with the calibration level?

\subsection{The calibration ladder does not move the command domain}

Table~\ref{tab:cmd} answers the second half first. Median commanded surge
in simulation is \SI{0.1428}{} to \SI{0.1438}{\metre\per\second} at every
level and in both geometries -- a spread of one part in a hundred and
forty across the entire ladder. The physical median is
\SI{0.123}{\metre\per\second} in \textsf{K0}. The absolute difference is
therefore \SIrange{0.0198}{0.0207}{\metre\per\second} at S0, S1, S2 and
S3 alike, and it does not shrink as the simulator becomes more realistic.

This is the expected result once stated, and it is worth stating because
it bounds what a boundary-level comparison can ever show. S3 acts
downstream of this signal by design -- that is what makes it a plant
model and not a controller edit -- so it cannot move the commanded surge
at all. S1 and S2 act upstream, but they change \emph{when} the planner
acts rather than how hard it pushes, and the planner's surge setpoints
are constants. A study that reported shrinking command-domain error
across such a ladder would be reporting something other than what it
claimed.

\begin{table}[t]
\centering
\caption{Command-domain comparison, committed planner. Simulated values
are medians over five frozen runs per cell \frozenmark; physical values
are medians over the admissible runs of each geometry \realmark
($n=4,2,2$). \emph{peak} is the largest lateral command: the planner
strafes at \SI{0.20}{\metre\per\second} and reaches
\SI{0.30}{\metre\per\second} only after entering its go-around phase.}
\label{tab:cmd}
\begin{tabular}{llrrrrr}
\toprule
domain & geo & $n$ & surge (\si{\metre\per\second}) & lateral frac.
& peak (\si{\metre\per\second}) & reversals \\
\midrule
sim S0 & K0 & 5 & 0.143 & 0.34 & 0.300 & 2 \\
sim S1 & K0 & 5 & 0.144 & 0.30 & 0.300 & 2 \\
sim S2 & K0 & 5 & 0.143 & 0.34 & 0.300 & 2 \\
sim S3 & K0 & 5 & 0.143 & 0.30 & 0.300 & 2 \\
\midrule
real & K0  & 4 & 0.123 & 0.42 & 0.200 & 1.5 \\
real & K1  & 2 & 0.114 & 0.51 & 0.200 & 1 \\
real & K1M & 2 & 0.115 & 0.50 & 0.200 & 2 \\
\bottomrule
\end{tabular}
\end{table}

\subsection{The manoeuvre that never happened}

The informative disagreement is not in a median but in a threshold. The
committed planner strafes at \SI{0.20}{\metre\per\second} while it
displaces to its clearing offset, and commands
\SI{0.30}{\metre\per\second} of sway with \SI{0.40}{\metre\per\second} of
surge only once it enters the go-around phase and begins running past
the obstacle. The peak lateral command therefore says, unambiguously and
without any spatial measurement, whether a run ever reached that phase.

Every one of the twenty frozen simulated runs in \textsf{K0} reached it,
at all four calibration levels. \emph{None of the eight physical runs
did.} All eight peak at exactly \SI{0.200}{\metre\per\second} of lateral
command and exactly \SI{0.150}{\metre\per\second} of surge -- the strafe
setpoint and the nominal surge, never the go-around values.

The physical vehicle, in other words, spent every run in the first phase
of a manoeuvre it never completed. That is precisely what the geometry
of Section~\ref{sec:deployment} predicts: the clearing offset is
\SI{2.5}{\metre} in a basin \SIrange{2.7}{3.0}{\metre} wide, so the
displacement condition that ends the strafe phase can never be satisfied,
and the planner strafes until the run ends or the operator intervenes.
The higher physical lateral fraction (\numrange{0.42}{0.51} against
\numrange{0.30}{0.34} in simulation) is the same fact seen from another
side: more of the run is spent strafing because the strafe never
terminates.

We regard this as the most useful sim-real finding in the study, and note
what made it available. It required no ground truth, no tracking and no
spatial reconstruction -- only the recognition that a controller's own
setpoint constants make its internal phase observable in the command
stream. Where external reference is unavailable, as it is on most small
vehicles most of the time, instrumenting the controller's phase
boundaries in this way recovers a genuine behavioural comparison from
data that otherwise supports only summary statistics.

\subsection{Commitment: measurable in one domain, predetermined in the
other}

Commitment distance is absent from the physical results for a reason
beyond missing instrumentation. The vehicle was inside engagement range
at release, and two of the four \textsf{K0} runs are already commanding
laterally at their first sample, so even with perfect tracking the
physical commitment distance would be a property of where the vehicle was
put rather than of what the planner decided. A sim-real error on that
quantity would compare a measurement against a geometric certainty.
Commitment \emph{time} remains meaningful and is given in
Table~\ref{tab:real}, censored where the run began already committed.
